% Portfolio rules in the GREU financial block.
% Describes the position rules for each sector and instrument.
% Explains the accounting concepts the rules depend on.
% Compile with lualatex.
\documentclass[11pt,a4paper]{article}

\usepackage{fontspec}
\setmainfont{Latin Modern Roman}
\setmonofont{DejaVu Sans Mono}[Scale=MatchLowercase]

\usepackage[margin=2.2cm]{geometry}
\usepackage{booktabs}
\usepackage{tabularx}
\usepackage{fancyvrb}
\usepackage{parskip}
\usepackage{amsmath}
\usepackage[hidelinks]{hyperref}

% Long code identifiers in running text cannot hyphenate. Give the paragraph
% builder a little slack rather than let a line run into the margin.
\emergencystretch=2em

\newenvironment{code}[1]
  {\par\vspace{0.4em}\noindent\texttt{\footnotesize #1}\par\vspace{-0.2em}
   \VerbatimEnvironment\begin{Verbatim}[fontsize=\scriptsize,frame=leftline,framesep=6pt,xleftmargin=6pt]}
  {\end{Verbatim}\vspace{0.2em}}

\title{Portfolio rules in the GREU financial block}
\author{}
\date{18 September 2026}

\begin{document}
\maketitle

This note describes how the model decides what each sector holds and what each
sector owes. It covers households, corporations, government, and the rest of the
world.

Section~\ref{sec:concepts} explains the accounting the rules rest on. Skip it if
you already know the European System of Accounts (ESA) and the growth adjustment.
The rest of the note reports the rules as they stand in the code, and then states
what follows from them.

\section{The concepts}
\label{sec:concepts}

\subsection{Sectors, instruments, and sides}
\label{sec:sectors}

The model splits the economy into five \emph{institutional sectors}: households
(\texttt{:Hh}), non-financial corporations (\texttt{:NonFinCorp}), financial
corporations (\texttt{:FinCorp}), government (\texttt{:Gov}), and the rest of the
world (\texttt{:RoW}). The last one holds everything outside the country.

Financial claims come in two \emph{instruments}. \texttt{:Debt} is lending and
borrowing. \texttt{:Equity} is ownership, mostly shares. Each sector can be on
either \emph{side} of an instrument. \texttt{:Assets} means it owns the claim,
\texttt{:Liab} means it owes it.

Five sectors, two instruments, two sides gives up to twenty stock cells per year.
The variable \texttt{vFinPosition\_s\_f[s,f,al,t]} holds them, and a mask keeps
only the cells the source reports in the calibration year.

Every claim has two ends. If a firm issues a share, some sector owns it. So for
each instrument, the total across sectors on the asset side must equal the total on
the liability side. This is the accounting fact behind the market clearing
equations in Section~\ref{sec:row}.

\subsection{How a stock changes}
\label{sec:stockchange}

A financial stock can change for three separate reasons, and ESA keeps them apart.

A \emph{transaction} is a purchase or a sale. Someone acts. If a household buys
shares for 100, its equity assets rise by 100 through a transaction.

A \emph{revaluation} is a price change on what the sector already holds. Nobody
acts. If the shares a household already owns gain 5 percent, its equity assets rise
through a revaluation. This is the capital gain.

An \emph{other change in volume} covers reclassification, write-off, and similar
bookkeeping events. The model takes these as given from the source.

The closing stock is the opening stock plus all three:

\begin{code}{SectorAccounts.jl:244--246}
vFinTransactions_f[s=sector, f=fin_instrument, al=ass_liab, t=t1:T],
vFinPosition_s_f[s,f,al,t] == vFinPosition_s_f[s,f,al,t-1]/fv + vFinTransactions_f[s,f,al,t]
                            + vFinReval_s_f[s,f,al,t] + vOtherChangesInVolume_f[s,f,al,t]
\end{code}

The distinction matters for this note. A sector that never transacts in an
instrument can still see its stock move, through revaluation alone. Three sectors
do exactly that with equity.

\subsection{Growth adjusted units and \texttt{fv}}

The model does not carry money values in current prices. It divides them by a
trend, so that a variable on a balanced growth path is flat. The trend factors are
in \texttt{GrowthInflationAdjustment.jl}: real growth \texttt{gq = 0.01},
inflation \texttt{gp = 0.02}, and the value factor
\[
  f_v = (1+g_q)(1+g_p) = 1.01 \times 1.02 = 1.0302 .
\]

So a value that is constant in these units grows 3.02 percent a year in money
terms. When an equation needs last year's stock measured in this year's units, it
writes \texttt{position[t-1]/fv}. That is what the division by \texttt{fv} in the
identity above means. It is not a discount and not a depreciation.

This has one consequence used throughout the note. In a steady state every
adjusted variable is flat. If a stock is still drifting at the end of the horizon,
the model has not reached a steady state, whatever the level says.

\subsection{Net financial assets}

A sector's \emph{net financial assets} are its total financial assets less its
total financial liabilities, in \texttt{vNetFinAssets[s,t]}. It is the sector's net
financial position: positive means it is a net creditor, negative a net debtor.

Because every liability is somebody's asset, the sum of net financial assets over
all five sectors must be zero. The model does not impose that as an equation. It
checks it, which Section~\ref{sec:row} returns to.

Net financial assets build up over time from the sector's own budget. If a sector
spends less than it earns, it lends the difference, and its net position improves.

\subsection{What a ``rule'' is, and what a ``residual'' is}

The model is a \emph{square} system: it has as many equations as unknowns, and it
solves them together. Each unknown needs exactly one equation to pin it down.

For each stock cell, the identity in Section~\ref{sec:stockchange} already determines the
transaction, and the revaluation identity determines the revaluation. So one more
equation is needed to determine the \emph{position} itself. That equation is what
this note calls the \textbf{rule} for the cell. In a
\texttt{SquareModels.@block}, the variable named before the comma is the one the
equation determines.

A sector module therefore never sets a transaction directly. It sets the position,
and the transaction comes out of the identity. This is why a rule that freezes a
stock has to be written as \texttt{vFinTransactions\_f == 0}. That is how you say
``this position moves only by revaluation''.

One cell per sector usually gets a special rule: its own balance sheet identity,
net financial assets equals assets less liabilities. That cell is not chosen by
anything. It absorbs whatever the other cells and the sector's net financial
assets leave over. This note calls it the \textbf{residual}, or the sector's
\textbf{plug}. Which cell a sector plugs with is a modelling choice, and it carries
most of the economics.

\subsection{Unit root}
\label{sec:unitroot}

A \emph{unit root} means a shock never fades. Take a stock that moves only by
revaluation, at rate $r$, in adjusted units:
\[
  V_t = \frac{V_{t-1}}{f_v} + r\,\frac{V_{t-1}}{f_v}
      = \frac{1+r}{f_v}\,V_{t-1}.
\]
For the adjusted stock to hold still, the model needs $r = f_v - 1$. Then the
coefficient on $V_{t-1}$ is exactly one. Nothing pulls the stock back toward
anything. A one-off hit of any size stays forever, and it compounds: a smaller
stock earns a smaller revaluation, which keeps the stock smaller.

This is the property that the companion note
\texttt{wealth\_channel\_diagnosis.md} measures for household equity.

\section{The rules, all of them}

\begin{table}[h]
\centering
\small
\begin{tabularx}{\textwidth}{@{}lXXXX@{}}
\toprule
Sector & Debt assets & Debt liabilities & Equity assets & Equity liabilities \\
\midrule
\texttt{:Hh}
  & \textbf{residual}
  & partial adjustment toward a share of consumption
  & share of all issued equity
  & no rule, exogenous \\
\addlinespace
\texttt{:NonFinCorp}
  & share of operating expenses
  & share of the capital replacement value
  & share of own equity liabilities
  & \textbf{residual} \\
\addlinespace
\texttt{:FinCorp}
  & share of all debt liabilities
  & share of own equity liabilities
  & zero transactions
  & \textbf{residual} \\
\addlinespace
\texttt{:Gov}
  & zero transactions
  & \textbf{residual}
  & zero transactions, both sides
  & zero transactions, both sides \\
\addlinespace
\texttt{:RoW}
  & market clearing
  & share of FinCorp debt assets
  & market clearing
  & share of domestic equity assets \\
\bottomrule
\end{tabularx}
\end{table}

Each sector's net financial assets come from its own budget identity plus
accumulation, independently of these rules.

\section{Households}

The budget identity gives net financial transactions, and accumulation gives net
financial assets. Within that total, debt liabilities follow a \emph{partial
adjustment} rule. The stock does not jump to its target. It closes a fixed
fraction of the gap each year, here one fifth, since
\texttt{rHhDebtAdjustment = 0.2}. The target is a fixed share of consumption, which
stands in for borrowing capacity.

\begin{code}{Households.jl:78--81}
# Debt liabilities move part of the way to a fixed share of consumption.
vFinPosition_s_f[s=[:Hh], f=[:Debt], al=[:Liab], t=t1:T],
vFinPosition_s_f[s,f,al,t] == (1 - rHhDebtAdjustment[t]) * vFinPosition_s_f[s,f,al,t-1]/fv
                          + rHhDebtAdjustment[t] * rHhDebtLiabilities2Consumption[t] * vC[t]
\end{code}

That is the only rule in the whole financial block with its own dynamics. Every
other rule holds within the year.

Households plug with debt assets. Extra saving becomes lending, which the code
comment states plainly. Equity was frozen at zero transactions until the change
described in the diagnosis note.

\texttt{mHhReturn} needs care, because the name suggests more than the equation
does.

\begin{code}{Households.jl:88--90}
# Extra household saving is held in debt assets.
mHhReturn[t=t1:T],
mHhReturn[t] * vFinPosition_s_f[:Hh,:Debt,:Assets,t-1]/fv == vFinIncome_s_f[:Hh,:Debt,:Assets,t]
\end{code}

It \textbf{reports} a number: interest received divided by the opening stock of
debt assets. It does not allocate anything. It enters the consumption Euler
equation, which is the condition that sets how much a household saves by comparing
utility today against utility tomorrow. So this return affects how much households
save, never what they hold it in.

\section{Corporations}

Both corporate sectors work the same way. Three cells follow fixed ratios, and
equity liabilities are the residual.

The ratios differ by what each sector is for. Non-financial corporations tie debt
liabilities to the replacement value of their capital, which is what a firm can
borrow against, and debt assets to operating expenses, which is working cash.
Financial corporations tie debt assets to all debt liabilities in the economy,
because lending to everyone else is their business, and their own debt liabilities
to their own equity liabilities, which is a \emph{leverage} ratio: how much a bank
borrows for each unit of its own capital.

The consequence is that \textbf{firm value is not a price, it is a plug}. Equity
liabilities absorb whatever the rest of the balance sheet leaves over. The present
value equation in \texttt{FinancialRevaluations.jl}, lines 85 to 88, does not set
the level of firm value. It sets the revaluation rate, and the equity issuance flow
absorbs the difference.

This shows up as a contradiction that is worth naming. \emph{Tobin's q} is the
market value of a firm divided by the replacement cost of its capital. When $q$
rises, capital is cheap next to what the market pays for it, so a firm should
invest more. In the productivity shock runs, investment rose while firm value fell
against the capital stock, so $q$ fell. Those two cannot both be right.

\section{Government}

Government differs from the other sectors in two ways.

First, it does not build its net financial transactions from a detailed budget
identity. It takes the net lending or borrowing balance directly. The
\emph{primary} balance is revenue less expenditure with interest left out. Adding
interest received and subtracting interest paid turns it into the full balance.

\begin{code}{Government.jl:102--105}
vGovBalance[t=t1:T],
vGovBalance[t] == vGovPrimaryBalance[t] + vFinIncome_s_f[:Gov,:Debt,:Assets,t] - vFinIncome_s_f[:Gov,:Debt,:Liab,t]

vNetFinTransactions[s=[:Gov], t=t1:T], vNetFinTransactions[s,t] == vGovBalance[t]
\end{code}

Second, it has the least portfolio freedom of any sector. Three of its four cells
are frozen. The equity rule covers both sides in one equation, because it is
written with \texttt{al=ass\_liab} instead of a single side.

\begin{code}{Government.jl:113--123}
# Gov neither buys nor sells equity; existing equity stocks follow non-transaction changes.
vFinPosition_s_f[s=[:Gov], f=[:Equity], al=ass_liab, t=t1:T], vFinTransactions_f[s,f,al,t] == 0

# Gov does not buy or sell debt assets; the stock follows non-transaction changes.
vFinPosition_s_f[s=[:Gov], f=[:Debt], al=[:Assets], t=t1:T], vFinTransactions_f[s,f,al,t] == 0

# Gov debt liabilities are residual given net financial assets.
vFinPosition_s_f[s=[:Gov], f=[:Debt], al=[:Liab], t=t1:T],
vNetFinAssets[s,t] == ∑(vFinPosition_s_f[s,f,:Assets,t] for f in fin_instrument)
                    - ∑(vFinPosition_s_f[s,f,:Liab,t] for f in fin_instrument)
\end{code}

Government plugs with debt liabilities. That is the right economics: a government
finances a deficit by issuing debt. It also means that holding net financial assets
fixed in a scenario holds gross debt close to fixed as well, because the asset side
cannot respond. A closure on net worth is therefore nearly a closure on gross debt.

Two consequences follow for scenario work.

The interest feedback is live. Debt interest enters the balance above, and the
balance drives the debt stock through accumulation. More debt raises interest,
which worsens the balance, which raises debt again. The core has no tax or spending
rule to break that loop, so a shock that moves the primary balance puts debt on a
permanent drift and no stock settles. Each scenario needs its own fiscal closure.
The productivity shock runs used one: hold \texttt{vNetFinAssets[:Gov]} on its
baseline path and let \texttt{tHhIncome} carry the adjustment.

Government equity has the same unit root as the other frozen equity cells. The
position is small, and a closure that fixes net financial assets makes the debt
side absorb it, so it does not drive results. It is listed for completeness.

\section{Rest of the world}
\label{sec:row}

The rest of the world is the residual counterparty in both markets. Two of its
four cells are \emph{market clearing} conditions. Market clearing here is not a
behavioural statement about prices. It is the accounting fact from Section~\ref{sec:sectors}:
every claim has a holder, so total assets must equal total liabilities for each
instrument. The rest of the world takes whatever the domestic sectors do not.

\begin{code}{RestOfWorld.jl:52--67}
# Debt assets clear the debt market. RoW is residual lender.
vFinPosition_s_f[s=[:RoW], f=[:Debt], al=[:Assets], t=t1:T],
vFinPosition_s_f[s,f,al,t] == ∑(vFinPosition_s_f[s2,:Debt,:Liab,t] for s2 in sector)
                          - ∑(vFinPosition_s_f[s2,:Debt,:Assets,t] for s2 in sector if s2 != :RoW)

# Debt liabilities are a fixed share of FinCorp debt assets.
vFinPosition_s_f[s=[:RoW], f=[:Debt], al=[:Liab], t=t1:T],
vFinPosition_s_f[s,f,al,t] == rForeignDebt[t] * vFinPosition_s_f[:FinCorp,:Debt,:Assets,t]

# Equity liabilities are a fixed share of domestic equity assets.
vFinPosition_s_f[s=[:RoW], f=[:Equity], al=[:Liab], t=t1:T],
vFinPosition_s_f[s,f,al,t] ==
  rForeignEquity[t] * ∑(vFinPosition_s_f[s2,:Equity,:Assets,t] for s2 in sector if s2 != :RoW)

# RoW equity assets clear the equity market.
vFinPosition_s_f[s=[:RoW], f=[:Equity], al=[:Assets], t=t1:T],
∑(vFinPosition_s_f[s2,:Equity,:Assets,t] for s2 in sector) ==
  ∑(vFinPosition_s_f[s2,:Equity,:Liab,t] for s2 in sector)
\end{code}

Read the two ratio rules as the foreign share of each portfolio.
\texttt{rForeignDebt} is the part of bank lending that goes abroad.
\texttt{rForeignEquity} is the foreign part of resident equity portfolios.

Now note what keeps the block square. Every other sector uses its own balance sheet
identity as the rule for one cell. The rest of the world does not. Its four
positions come from the four rules above, and its net financial assets come from
its budget identity, all independently. So for this one sector the identity is not
imposed. It is only checked:

\begin{code}{SectorAccounts.jl:277--278}
@test_constraint("Summing vNetFinAssets over sectors"; atol=30.0, rtol=1e-6)
vNetFinAssets[s=[:Hh], t=t1:T], ∑(vNetFinAssets[s,t] for s in sector) == 0.0
\end{code}

A \texttt{@test\_constraint} is not solved. It is verified after the solve, within
the tolerance given. This one is the only check that the financial block adds up
across sectors, so it is not decoration. In the productivity shock runs the sum came
out near $-1$ MEUR against a tolerance of 30, so it holds. The identity for the
rest of the world is then implied by the other four sectors plus the two clearing
conditions, rather than assumed.

\section{What this means}

\textbf{There is no portfolio choice anywhere in the model}, in the economic sense.
No sector compares the return on debt against the return on equity and shifts
between them. Every allocation is a fixed ratio, a frozen stock, or an accounting
residual. For a model built on national accounts that is a defensible starting
point, but it should be said out loud, because the variable names suggest
otherwise.

\textbf{The choice of plug carries the economics.} Households plug with debt
assets, so extra saving becomes lending. Government plugs with debt liabilities, so
a deficit becomes debt. Both are sound. Both corporate sectors plug with equity
liabilities, so anything the balance sheet leaves over lands on their market value.
That is the one we would question, because it makes firm value an accounting
residual instead of a price.

\textbf{Frozen equity brings a unit root.} Equity sits under a zero transaction
rule for households before the change, for financial corporations, and for
government on both sides. By Section~\ref{sec:unitroot} each of those stocks has a unit root, so
any hit to its revaluation rate is permanent. Consumption in this model depends on
the level of household financial wealth, so that permanence reaches the real
economy. The diagnosis note gives the measurements.

\textbf{One loop deserves attention when reading a shock.} Household equity assets
are part of domestic equity assets. Domestic equity assets form the base for
\texttt{rForeignEquity}, which sets the equity liability of the rest of the world.
That liability is part of total issued equity, which is what the household share
rule reads. So the rule feeds back into its own right hand side within the same
year. The strength of that feedback, its \emph{loop gain}, is near
\texttt{rHhEquityAssets2TotalEquity} times \texttt{rForeignEquity}. Two small
shares multiplied give a small number, so the loop is weak, but it is not zero, and
it sits in the same year as the two clearing conditions.

\end{document}
